\section{Results}\label{sec:results}
Table~\ref{tab:hyp} lists the primary hypotheses of studies 2 to 7, as fixed in their protocols, with the observed
differences and outcomes. Each row compares two methods on the score named in the fourth column, and a positive
difference favours the first method. The subsections follow the order in which the studies were designed, so that each
design decision can be traced to the result that motivated it.

\inputtable{tab_hyp}

\subsection{Planning model and control of computation}\label{sec:study1}\label{sec:study2}
Study~1 made only the stopping rule execution-aware and kept the nominal deadlines. AnytimeSwarm-R needed
\val{JMsRedAS}\% less computation per call than AnytimeSwarm, but its advantage over ALNS with 40\,ms was not
significant after the Holm correction (Supplementary Section~\ref{S-sec:Sstudy1}). Study~2 therefore moved the error
model from the stopping rule into the planning model. \mname{} without fallback reached the highest utility,
\val{KUProp}\%. Its advantage over the six baselines of H1 ranged from \val{KDmBuf} to \val{KDmForty}\,pp, and over the
conventional planner it was \val{KDConv}\,pp, or \val{KURelConv}\% (H2). With the effective deadlines, the stopping
rule was equivalent to a fixed budget of 40\,ms (H3), and its value lay mainly in avoiding the late plans of long
budgets on slow planners (Supplementary Section~\ref{S-sec:Sstudy2}). Without fallback, however, \mname{} was too
cautious under high disturbances with halved windows, where it served \val{KDBufWHalfHigh}\,pp relative to the constant
buffer, because many effective deadlines fell before the earliest possible arrival. This result led to the fallback.

\subsection{Margin shape and fallback}\label{sec:study3}\label{sec:study4}
Study~3 tested the nominal fallback on new seeds. It added \val{LDPropAll}\,pp to \mname{} without fallback (H5) and
\val{LDBufHighWHalf}\,pp over the constant buffer in the regime where study~2 had lost (H6), at the price of more tasks
served late, since a fallback task is planned against its true deadline (Supplementary Section~\ref{S-sec:Sstudy3}).

Table~\ref{tab:s4} compares the planning models of study~4 on a third set of seeds. All models use the same search for
40\,ms, so they differ only in the margin and the fallback. The columns give the utility over all configurations and
per regime, the tasks served late per mission and the search iterations per millisecond. Three patterns stand out.

First, a margin that grows with the lead time is better than none or a constant one. \mname{} served
\val{XDabsBufAll}\,pp more reward than the constant buffer (95\% CI \val{XCIBufAll}, H12) and \val{XDabsFortyAll}\,pp
more than the nominal model (\val{XCIFortyAll}, H13), that is, \val{QURelCCBuf}\% and \val{QURelCCForty}\% more. The
advantage held in each of the four regimes, and with repair only and no search it was still \val{XDabsRepCCRep}\,pp.
Over the three independent sets of seeds of studies 2 to 4, every advantage over the constant buffer, the nominal model
and the conventional planner was significant (Supplementary Fig.~\ref{S-fig:replication}).

Second, the three margins that grow with the lead time reached utilities within \val{XRangeScaled}\,pp of each other.
The proportional buffer served \val{XDabsPropBufAll}\,pp more than \mname{} (H9) and moment propagation
\val{XDabsMPAll}\,pp more (H10), so the calibration adds little utility beyond a well-tuned slope. The closed form was
the cheapest of the three, with \val{XDabsItMsMP} more search iterations per millisecond than moment propagation.

Third, every chance-constrained model needs the fallback. Without it, moment propagation reached the lowest utility of
all models, \val{QUMPnf}\%, \val{XDabsFbMPAll}\,pp below its version with fallback (H11), while it served the fewest
tasks late. A hard chance constraint on every task, however it is evaluated, thus gives up tasks that are still worth
attempting. The error model also ranked the served tasks by risk. Under high disturbances, \val{QCovObsCCHigh}\% of
them were served on time against a predicted \val{QCovPredCCHigh}\% (Supplementary Fig.~\ref{S-fig:Scoverage}).

The online variant, which does not know the disturbance level, served \val{XDabsOnlCC}\,pp less than the model
fitted for the known level and still well above the nominal model (Table~\ref{tab:s4}). This comparison was secondary,
so the protocol did not test non-inferiority, but it indicates that the model can be calibrated during the mission.

\inputtable{tabm_s4}

\subsection{Software-in-the-loop flights}\label{sec:s2-tierc}
In study~2, \val{KCRuns} held-out flights with mid disturbances and the original windows did not confirm the advantage
of \mname{} without fallback (H4, \val{KCDForty}\,pp). Tier B predicted the utility of single flights closely
($r={}$\val{KBCr}) but had expected only \val{KBgzDForty}\,pp on these missions (Supplementary
Section~\ref{S-sec:Ss2-tierc}). Study~4C was therefore placed in the regime with the largest predicted effect and flew
\val{GNFlights} flights on \val{HNB} new missions. All flights ran at a valid real-time factor, with a mean of
\val{GRTF}.

Fig.~\ref{fig:demo} shows one development mission flown with the graphical interface. Panels (a) and (b) show the
Gazebo scene, in which every task is a pillar coloured by its outcome, and panels (c) and (d) the trajectories recorded
by PX4. The nominal model sent a UAV to a task that it reached after the deadline, whereas \mname{} served the same
task on time with another UAV and then one more task. It served \val{ZDemoServedProp} of the \val{ZDemoNTasks} tasks on
time against \val{ZDemoServedNom}, with \val{ZDemoLateProp} instead of \val{ZDemoLateNom} late arrival. This pair of
flights illustrates the mechanism and is not part of any test.

\begin{figure*}[t]
\centering
\fig{fig_demo_gazebo}{0.80\linewidth}\\[2pt]
\fig{fig_demo_tracks}{0.54\linewidth}
\caption{Illustrative software-in-the-loop flight on a development mission. (a), (b) Gazebo scene. (c) \mname{} with
nominal fallback and (d) nominal model: trajectories recorded by PX4.}\label{fig:demo}
\end{figure*}

Fig.~\ref{fig:s4c} summarises the held-out flights of study~4C. For each primary comparison, it shows the tier C
result next to the tier B result on the same missions and the tier B estimate on development seeds that sized the
study. The fallback carried over to the flights. With it, \mname{} served \val{XDabsTcNoFb}\,pp more reward than
without it (95\% CI \val{XCITcNoFb}, H15), \val{GUFB}\% against \val{GUNoFb}\%, or \val{HURelNoFb}\% more, and it was
better in \val{HWinNoFb} of the \val{GN} missions. Against the nominal model, it served \val{XDabsTcNom}\,pp less
(\val{XCITcNom}), so H14 was not supported. Tier B gave the same picture on the same missions, \val{XDabsTcBNom}\,pp
less than the nominal model and \val{XDabsTcBNoFb}\,pp more than without fallback, and it predicted single flights
with a correlation of \val{HBCr}. The gap to the development estimate therefore came from the sampled missions rather
than from the transfer to the autopilot.

\begin{figure}[t]
\centering
\fig{fig_s4c}{\linewidth}
\caption{Study 4C: utility of \mname{} with nominal fallback minus the nominal model (H14) and minus \mname{} without
fallback (H15), in tier C, in tier B on the same missions and in tier B on development seeds. Error bars are 95\%
intervals.}\label{fig:s4c}
\end{figure}

\subsection{Cost of a late arrival}\label{sec:study5}
The utility counts only the tasks served on time, so a late arrival earns nothing and costs nothing. Re-scored with a
late cost, the earlier runs showed that the nominal fallback pays off only when a late arrival costs less than about
one task reward. Its advantage vanished at a late cost of \val{PSThreeFbPropStar} task rewards in study~3,
\val{PSFourMPMPnfStar} for moment propagation in study~4 and \val{PCFbNoFbStar} in tier C (Supplementary
Fig.~\ref{S-fig:Slatepen}). Study~5 tested the cost-aware fallback on new seeds.

Table~\ref{tab:s5} gives the mean score of each method for four late costs and, in the last column, the tasks served
late at one task reward. The cost-aware fallback and the tuned competitors are run for the late cost of each column,
so each column is a separate comparison. At one task reward, the cost-aware fallback scored \val{FUCAOneOne}\% against
\val{FUCCfbOne}\% for the nominal fallback, \val{XDabsFiveCCfbOne}\,pp more (95\% CI \val{XCIFiveCCfbOne}, H16), and
it served \val{FLateRedCCfb}\% fewer tasks late. It was also \val{XDabsFivePropOne}\,pp above the proportional buffer
tuned for this cost (H17) and \val{XDabsFiveCCnfOne}\,pp above \mname{} without fallback with a tuned safety factor
(H18). Both competitors had been selected from sixteen and five candidates for this cost, whereas the cost-aware
fallback kept the safety factor of the utility and had no cost-specific tuning.

\inputtable{tabm_s5}

Fig.~\ref{fig:s5} shows how these differences change with the late cost. The advantage over the nominal fallback grows
with the cost, because the nominal fallback attempts every reachable task whatever a late arrival costs. Against the
tuned competitors, the cost-aware fallback was on par with the proportional buffer and ahead of the tuned safety factor
at half a task reward, and ahead of both at one task reward. At two task rewards, the proportional buffer tuned for
this cost was better by \val{XDabsFivePropTwo}\,pp. A fallback task then needs a probability of at least
\val{FThetaTwo} of being on time, and the delay that it causes to later tasks on its route, which the rule ignores,
weighs more. The competitors, however, have to be tuned again for every cost. Tuned for one task reward, the
proportional buffer stayed \val{XDabsFivePropFixTwo}\,pp behind the cost-aware fallback at two task rewards.

\begin{figure}[t]
\centering
\fig{fig_s5}{\linewidth}
\caption{Study 5: score of \mname{} with the cost-aware fallback minus each method against the late cost. Error bars
are 95\% intervals over 270 missions.}\label{fig:s5}
\end{figure}

In practical units, per 100 tasks the cost-aware fallback served \val{YOnTimeCA} on time and \val{YLateHundredCA}
late, the nominal fallback \val{YOnTimeCCfb} and \val{YLateHundredCCfb}, and the nominal model \val{YOnTimeForty} and
\val{YLateHundredForty}. The fleet flew \val{YKmCA}\,km per mission instead of \val{YKmCCfb} and \val{YKmForty}\,km,
consistent with fewer flights to tasks that were likely to be reached too late.

\subsection{Comparison with CBBA and moment propagation}\label{sec:study6}
Table~\ref{tab:s6} adds CBBA and moment propagation with the cost-aware fallback. All methods with this fallback were
run at one task reward and are scored at every late cost from the same runs, and the last two columns give the late
tasks per mission and the planning time per call. With the score $u_1$, \mname{} exceeded CBBA by
\val{WDabsCBBAOne}\,pp (95\% CI \val{WCICBBAOne}, H20), and on utility alone by \val{WDabsCBBAZero}\,pp, or
\val{WURelCACBBA}\%. The difference lies in the planning model rather than in the search. CBBA converged after
\val{WItersCBBA} consensus rounds and \val{WMsCBBA}\,ms per call and reached almost the utility of ALNS on the nominal
model. Planned on the \mname{} deadlines, it gained \val{WDabsCBBAModelOne}\,pp (H21) and came within
\val{WDabsCACBBACAOne}\,pp of \mname{} with \val{WMsCBBACA}\,ms per call. The effective deadlines can thus be handed
to an auction-based allocator without changing its bidding rule. The advantage of the cost-aware over the nominal
fallback within \mname{}, \val{WDabsCACCfbOne}\,pp, replicated H16 on a second set of seeds.

\inputtable{tabm_s6}

Fig.~\ref{fig:s6} plots the primary differences of study~6 against the late cost, the comparisons with CBBA in panel
(a) and the fallbacks within moment propagation in panel (b). Within moment propagation, the
cost-aware fallback scored \val{WDabsMPFbOne}\,pp above the nominal fallback (H22) and \val{WDabsMPNoFbOne}\,pp above
no fallback (H23) at one task reward. As in the closed form, it gave up some utility when lateness was free, broke even
at half a task reward and gained most at two, so the rule is not tied to the closed form. With the cost-aware fallback,
moment propagation scored \val{WDabsCAMPCAOne}\,pp above the closed form at one task reward and
\val{WDabsCAMPCATwo}\,pp at two, whereas the closed form was \val{WDabsCAMPCAZero}\,pp ahead when lateness was free.
Once lateness has a cost, the variance that moment propagation tracks leg by leg pays off.

\begin{figure}[t]
\centering
\fig{fig_s6}{\linewidth}
\caption{Study 6: score differences against the late cost. (a) Comparisons with CBBA. (b) Cost-aware fallback within
moment propagation. Error bars are 95\% intervals over 270 missions.}\label{fig:s6}
\end{figure}

\subsection{Late-cost model and errors outside the model}\label{sec:study7}
Fig.~\ref{fig:s7}a compares the two cost models of the rule with the proportional buffer tuned for the same score and
late cost. Each row is one score and late cost, and a point right of zero favours the cost-aware fallback. At one task
reward, matching the rule to the score changed little. With $u_1$, the fixed-cost rule minus the proportional rule
was \val{UDHTwentyFour}\,pp, not a significant difference (H24), and the fixed-cost rule minus the tuned buffer was
\val{UDHTwentyFive}\,pp (H25),
which replicates H17 on new seeds. With $u^w_1$, the proportional rule was non-inferior to the buffer tuned for that score
(\val{UDHTwentySix}\,pp, H26). At twice the cost, each rule did better on its own score. With $u_2$, the fixed-cost
rule no longer differed significantly from the tuned buffer, and with $u^w_2$ both rules stayed behind it. The score
$u_\lambda$ charges a late arrival the realised mean reward of the mission, whereas the fixed-cost rule assumes the mean
reward $\bar w=2.1$ of the generator, because the total reward is not known while tasks are still arriving. Study~7a was
therefore also scored with $u'_\lambda=U-\lambda\bar w\,n_{\mathrm{late}}/R$, which charges exactly the cost that the
rule assumes, and the conclusions did not change (Supplementary Table~\ref{S-tab:Ss7const}).

\begin{figure*}[t]
\centering
\fig{fig_s7}{0.92\linewidth}
\caption{Study 7. (a) Cost-aware fallback minus the tuned proportional buffer for each score and late cost; open
markers have an interval that includes zero. (b) Primary comparisons under the calibrated errors and four departures
from the error model; dashed lines mark the control values. Error bars are 95\% intervals over 270 missions.}\label{fig:s7}
\end{figure*}

Fig.~\ref{fig:s7}b shows the two primary comparisons of study~7b under the calibrated errors, shaded as control, and
under four departures from the error model. The advantages hardly changed. With the nominal fallback, \mname{} served
\val{UDabsHTwentySevenBias} to \val{UDabsHTwentySevenOutlier}\,pp more than the nominal model (H27), against
\val{UBCtlCCfb}\,pp under the calibrated errors, and the cost-aware fallback scored \val{UDabsHTwentyEightPareto} to
\val{UDabsHTwentyEightOutlier}\,pp more than the nominal fallback (H28). All eight comparisons were significant. The
smallest advantage occurred with a speed bias per UAV, which correlates the errors along a route. With rare long
delays, the predicted probability became slightly optimistic in the upper tail, with \val{VCalObsHiSOutlier}\% of the
tasks served on time against a predicted \val{VCalPredHiSOutlier}\% (Supplementary Section~\ref{S-sec:Scalib}).
